% Generated draft: requires mathematical and source review.
\documentclass[11pt]{article}
\usepackage[margin=25mm]{geometry}
\usepackage{fontspec}
\setmainfont{Latin Modern Roman}
\usepackage{amsmath,amssymb,mathtools,unicode-math}
\setmathfont{Latin Modern Math}
\usepackage{enumitem,needspace,xurl,hyperref}
\setlength{\parindent}{0pt}
\setlength{\parskip}{5pt}
\setlength{\emergencystretch}{3em}
\newcommand{\N}{\mathbb{N}}
\newcommand{\Z}{\mathbb{Z}}
\newcommand{\Q}{\mathbb{Q}}
\newcommand{\R}{\mathbb{R}}
\newcommand{\C}{\mathbb{C}}
\newcommand{\F}{\mathbb{F}}
\newcommand{\A}{\mathbb{A}}
\newcommand{\PP}{\mathbb{P}}
\newcommand{\E}{\mathbb{E}}
\newcommand{\eps}{\varepsilon}
\newcommand{\dd}{\,\mathrm{d}}
\begin{document}
\section*{UnsolvedMath ID: 11000060}
\textbf{Source problem code:} catalogue entry
\subsection{Definitions and mathematical statement}
Let \(\mathcal I_g\) and \(\mathcal K_g\) be the source's Torelli group and subgroup \(\mathcal K_g\). In the usual closed oriented surface setup \(\mathcal I_g\) acts trivially on integral first homology, while \(\mathcal K_g\) denotes the Johnson kernel, conventionally generated by separating twists in the appropriate genus range. The excerpt does not supply that definition, its genus hypotheses, or its surface convention.

\emph{Problem 5.15 (explicit cycles).} Explicitly construct infinitely many linearly independent rational homology classes represented by cycles in each of
\[
 H_*(\mathcal I_g;\Q)\qquad\text{and}\qquad H_*(\mathcal K_g;\Q).
\]
For instance, in an explicitly specified chain model, the requested data are cycles \(z_i\) with \(\partial z_i=0\) such that for every finite rational combination,
\[
 \sum_{i=1}^{r}q_i[z_i]=0\quad\Longrightarrow\quad q_1=\cdots=q_r=0.
\]
This requires independence of homology classes, not merely independence of chains. Whether genus is fixed, which genera are intended, and whether all cycles must lie in one degree are unstated; the homological reading follows the word ``cycles'' and the torsion-bearing degree-one context below. The source emphasizes the difficulty of finding even explicit nonzero cycles.

The following discussion reports Johnson's calculation, citing [Jo1] as a summary,
\[
 H_1(\mathcal I_g;\Z)\cong
 \bigl(\textstyle\bigwedge^3 H\bigr)/H\ \oplus B_2,
 \tag{16}
\]
where the source says \(B_2\) consists of \(2\)-torsion. The usual interpretation is \(H=H_1(S_g;\Z)\), with the quotient using the symplectic embedding of \(H\) in \(\bigwedge^3H\), but the excerpt does not define this embedding, \(B_2\), or the valid genus range. In particular, this cannot be a literal \(H^1(-;\Z)\) decomposition with torsion. The source contrasts an algebraic exterior-power part with the topological \(B_2\) contribution from the Rochlin invariant [BC, Jo1]; the torsion contribution vanishes with rational coefficients and so does not itself furnish the requested rational families. It also makes an analogy with \(H_1\) of the Torelli group for \(\operatorname{Out}(F_n)\), saying an algebraic part occurs there but the Rochlin part does not. This is retained as a qualitative source comparison, not an asserted identical module formula.
\subsection{Short English statement}
Construct explicit infinite families of independent rational homology classes for both the Torelli group and \(\mathcal K_g\). The source's integral first-homology formula and Rochlin-derived torsion are context; neither fixes the degrees of the desired families nor supplies rational classes from the torsion term.
\subsection{Sources}
\begin{itemize}
\item Local UnsolvedMath registry, permanent ID 11000060, source code catalogue entry, Problem 5.15, ``Explicit cycles,'' PDF page 48. No individual proposer is supplied.
\item Farb, \emph{Problems on Mapping Class Groups and Related Topics}. Supplied URL: \url{https://www.math.uchicago.edu/~farb/papers/mcgbook.pdf}; recorded access 2026-07-29; extraction pdf-font-structure. The edited volume contains 329 explicitly typeset headings across chapters, with repeated local numbering disambiguated by page.
\item Johnson [Jo1] is cited for the summary of the integral degree-one calculation, and [BC, Jo1] for its Rochlin contribution; full bibliographic entries are absent. The formula is numbered (16) in the supplied text.
\item Supplied status open and OPEN-TRIAGE. The summary says only ``Cannot verify without statement'' and gives no actual assessment of the provided cycle question. Current primary-literature status and redistribution terms for the publicly accessible author/publisher copy remain NEEDS\_REVIEW.
\end{itemize}
\end{document}
